% !TeX program = pdflatex
\documentclass[11pt,a4paper]{amsart}

\newcommand{\notelanguage}{finnish}
% Shared preamble for course notes and exercise sets.

\usepackage[T1]{fontenc}
\usepackage{lmodern}
\usepackage[english]{babel}

\usepackage[a4paper,margin=30mm]{geometry}
\usepackage{microtype}
\usepackage{graphicx}
\usepackage{enumitem}

\usepackage{amsmath}
\usepackage{amssymb}
\usepackage{amsthm}
\usepackage{mathtools}
\usepackage{aliascnt}

\usepackage[hidelinks,pdfusetitle]{hyperref}

\numberwithin{equation}{section}
\setlist{itemsep=0.25em,topsep=0.5em}

\theoremstyle{plain}
\newtheorem{theorem}{Theorem}[section]
\newaliascnt{lemma}{theorem}
\newtheorem{lemma}[lemma]{Lemma}
\aliascntresetthe{lemma}
\newaliascnt{proposition}{theorem}
\newtheorem{proposition}[proposition]{Proposition}
\aliascntresetthe{proposition}
\newaliascnt{corollary}{theorem}
\newtheorem{corollary}[corollary]{Corollary}
\aliascntresetthe{corollary}

\theoremstyle{definition}
\newaliascnt{definition}{theorem}
\newtheorem{definition}[definition]{Definition}
\aliascntresetthe{definition}
\newaliascnt{example}{theorem}
\newtheorem{example}[example]{Example}
\aliascntresetthe{example}
\newaliascnt{problem}{theorem}
\newtheorem{problem}[problem]{Problem}
\aliascntresetthe{problem}

\theoremstyle{remark}
\newaliascnt{remark}{theorem}
\newtheorem{remark}[remark]{Remark}
\aliascntresetthe{remark}

\usepackage[nameinlink,noabbrev]{cleveref}

\crefname{theorem}{theorem}{theorems}
\crefname{lemma}{lemma}{lemmas}
\crefname{proposition}{proposition}{propositions}
\crefname{corollary}{corollary}{corollaries}
\crefname{definition}{definition}{definitions}
\crefname{example}{example}{examples}
\crefname{problem}{problem}{problems}
\crefname{remark}{remark}{remarks}

\DeclareMathOperator{\im}{im}
\DeclareMathOperator{\rank}{rank}
\DeclarePairedDelimiter{\abs}{\lvert}{\rvert}
\DeclarePairedDelimiter{\norm}{\lVert}{\rVert}
\DeclarePairedDelimiter{\set}{\lbrace}{\rbrace}

\usepackage{tikz}

\title{Analyysi I\\Harjoituksen 3 ratkaisut}
\author{}
\date{}

\begin{document}

\exerciseheading{Analyysi I}{Harjoituksen 3 ratkaisut}

Tässä harjoituksessa merkitään \(\NNone = \{1,2,3,\ldots\}\).

\begin{exercise}
  Olkoon \((x_n)\) lukujono, jolle \(x_n = \frac{3n}{n + 37}\). Osoita suoraan
  määritelmästä, että lukujono suppenee kohti lukua \(3\). Osoita lisäksi,
  että lukujono on kasvava.
\end{exercise}

\begin{solution}
  Olkoon \(\epsilon > 0\). Valitaan \(N \in \NNone\) siten, että \(N >
  \frac{111}{\epsilon}\). Tällöin
  \[
    \abs{x_n - 3} = \frac{111}{n + 37} < \frac{111}{n} < \frac{111}{N}
    < \epsilon,
  \]
  kun \(n > N\). Täten \(\lim_{n \to \infty} x_n = 3\). Osoitetaan vielä,
  että \((x_n)\) on kasvava. Huomataan, että
  \[
    \frac{x_{n + 1}}{x_n} = \frac{\left(\frac{3(n + 1)}{(n + 1) +
    37}\right)}{\left(\frac{3n}{n + 37}\right)} = \frac{3n^2 + 114n +
    111}{3n^2 + 114n} = 1 + \frac{37}{n(n + 38)} > 1.
  \]
  Koska \(x_n > 0\) ja osamäärä \(\frac{x_{n + 1}}{x_n}\) on suurempi
  kuin \(1\),
  on voimassa \(x_n < x_{n + 1}\) kaikilla \(n \in \NNone\). Täten
  lukujono \((x_n)\) on aidosti kasvava ja siten myös kasvava.
\end{solution}

\begin{exercise}
  \leavevmode\par
\begin{enumerate}[label=\alph*), nosep]
  \item Olkoon \((x_n)\) suppeneva lukujono. Määritellään lukujono \((y_n)\)
    asettamalla, että \(y_n = \frac{1}{\sqrt{n}}x_n\). Osoita, että lukujono
    \((y_n)\) suppenee.
  \item Oletetaan, että lukujono \((x_n)\) on rajoitettu. Määritellään lukujono
    \((y_n)\) kuten (a)-kohdassa. Mitä voidaan sanoa lukujonon \((y_n)\)
    suppenemisesta?
\end{enumerate}
\end{exercise}

\begin{solution}
\leavevmode\par
\begin{enumerate}[label=\alph*), nosep]
\item Olkoon \(x = \lim_{n \to \infty} x_n\) ja \((z_n)\) lukujono, jolle
  \(z_n = \frac{1}{\sqrt{n}}\). Olkoon \(\epsilon > 0\). Valitaan \(N
  \in \NNone\)
  siten, että \(N > \frac{1}{\epsilon^2}\). Tällöin
  \[
    \abs{z_n - 0} = \frac{1}{\sqrt{n}} < \epsilon,
  \]
  kun \(n > N\), eli \(\lim_{n \to \infty} z_n = 0\). Lauseen (2.2.8) nojalla
  \[
    \lim_{n \to \infty} y_n = \lim_{n \to \infty} z_nx_n = \lim_{n
    \to \infty} z_n \cdot \lim_{n \to \infty} x_n = 0\cdot x = 0.
  \]
  Lukujono \((y_n)\) siis suppenee (kohti lukua \(0\)).

\item Koska \((x_n)\) on rajoitettu, on olemassa \(M \in \RR\), jolla
  \(\abs{x_n} \leq M\) kaikilla \(n \in \NNone\). Osoitetaan, että
  \(\lim_{n \to \infty} y_n = 0\). Olkoon \(\epsilon > 0\). Valitaan
  \(N \in \NNone\) siten, että \(N > \frac{M^2}{\epsilon^2}\). Tällöin
  \[
    \abs{y_n - 0} = \abs*{\frac{1}{\sqrt{n}}x_n} \leq
    M\frac{1}{\sqrt{n}} < \epsilon,
  \]
  kun \(n > N\). Lukujono \((y_n)\) suppenee siis kohti lukua \(0\).
\end{enumerate}
\end{solution}

\begin{exercise}
Anna esimerkki lukujonosta, joka
\begin{enumerate}[label=\alph*), nosep]
\item on kasvava, mutta ei aidosti kasvava, ja hajaantuu.
\item on aidosti vähenevä ja suppenee.
\end{enumerate}
\end{exercise}

\begin{solution}
\leavevmode\par
\begin{enumerate}[label=\alph*), nosep]
\item Olkoon \((x_n)\) lukujono, jolle
\[
x_1 = 2, \quad x_n = n, \quad n \geq 2.
\]
Huomataan, että \(x_1 = x_2\) ja \(x_n < x_{n + 1}\), kun \(n \geq 2\).
On siis voimassa \(x_n \leq x_{n + 1}\) kaikilla \(n \in \NNone\), eli
lukujono on kasvava, mutta ei aidosti kasvava. Olkoon \(M \in \RR\).
Valitaan \(N \in \NNone\) siten, että \(N > \maxset{2, M}\). Tällöin
\[
x_n = n > N > M
\]
kaikilla \(n > N\), joten lukujono kasvaa rajatta. Lemman 2.3.8 nojalla
lukujono \((x_n)\) hajaantuu.
\item Olkoon \((x_n)\) lukujono, jolle \(x_n = 1 + \frac{1}{n} =
\frac{n + 1}{n}\). Osoitetaan, että \(\lim_{n \to \infty} x_n = 1\).
Olkoon \(\epsilon > 0\). Valitaan \(N \in \NNone\) siten, että
\(N > \frac{1}{\epsilon}\). Tällöin
\[
\abs{x_n - 1} = \abs*{(1 + \frac{1}{n}) - 1} = \frac{1}{n} <
\frac{1}{N} < \epsilon,
\]
kun \(n > N\). Täten \(\lim_{n \to \infty} x_n = 1\). Osoitetaan vielä, että
\((x_n)\) on aidosti vähenevä. Huomataan, että
\[
x_{n + 1} = 1 + \frac{1}{n + 1} < 1 + \frac{1}{n} = x_n.
\]
Lukujono \((x_n)\) on siis aidosti vähenevä.
\end{enumerate}
\end{solution}

\begin{exercise}
Olkoon \(E\) reaalilukujen \(\RR\) epätyhjä alhaalta rajoitettu osajoukko.
Osoita, että jokaista \(\epsilon > 0\) kohti on olemassa sellainen \(x \in E\),
että \(\inf E \leq x < \inf E + \epsilon\).
\end{exercise}

\begin{solution}
Koska \(\inf E\) on joukon \(E\) alaraja, \(\inf E \leq x\) kaikilla
\(x \in E\).
Tehdään vastaoletus, että \(\inf{E} + \epsilon \leq x\) jollain
\(\epsilon > 0\) ja kaikilla \(x \in E\). Tällöin \(\inf{E} + \epsilon\)
on joukon \(E\) alaraja sekä \(\inf{E} + \epsilon > \inf{E}\).
Tämä on ristiriidassa sen kanssa, että \(\inf{E}\) on joukon \(E\) suurin
alaraja. Täten alkuperäisen väitteen on oltava tosi.
\end{solution}

\begin{exercise}
Tutustu kirjan Lauseeseen 2.3.4 (Bolzanon-Weierstrassin lause) ja sen
todistukseen. Bolzanon-Weierstrassin lauseen todistuksessa käytettiin joukkojen
\(\set{x_n \colon n \geq k}\) supremumeja. Muuta todistusta niin että se käyttää
samojen joukkojen infimumeja.
\end{exercise}

\begin{solution}
Olkoon \((x_n)\) rajoitettu lukujono ja \(A_k = \inf{\set{x_n \colon
n \geq k}}\). Koska \((x_n)\) on rajoitettu, niin \(A_k \in \RR\) kaikilla
\(k \in \NNone\). Koska
\(\set{x_n \colon n \geq k + 1} \subseteq \set{x_n \colon n \geq k}\),
pätee \(A_k \leq A_{k + 1}\). Jos \(m \leq x_n \leq M\) kaikilla
\(n \in \NNone\), niin myös \(m \leq A_k \leq M\) kaikilla
\(k \in \NNone\). Lukujono \((A_k)\) on siis kasvava ja rajoitettu,
joten se suppenee. Olkoon \(a = \lim_{k \to \infty} A_k\).

Valitaan osajono lukujonosta \((x_n)\). Tehtävän (4) mukaan lukujonosta
\((x_n)\) voidaan valita jäsen \(x_{n_{1}}\) siten, että
\[
x_{n_{1}} \leq A_1 + 1.
\]
Edetään induktiivisesti: oletetaan, että luvut \(x_{n_{1}}, \cdots, x_{n_{k}}\)
on valittu. Valitaan lukujonosta \((x_n)\) jäsen \(x_{n_{k + 1}}\),
jolle \(n_{k + 1} \geq n_{k} + 1\) ja
\[
x_{n_{k + 1}} \leq A_{n_k + 1} + \frac{1}{k + 1}.
\]
Koska \(n_k + 1 \to \infty\), kun \(k \to \infty\), ja \(A_k \to a\),
myös \(A_{n_k + 1} \to a\). Näin osajonolle \((x_{n_{k}})\) pätee
\[
A_{n_k + 1} \leq x_{n_{k + 1}} \leq A_{n_k + 1} + \frac{1}{k + 1}.
\]
Koska \(A_{n_k + 1} \to a\) ja \(\frac{1}{k + 1} \to 0\) kun \(k \to \infty\),
suppiloperiaatteen nojalla \(x_{n_{k}} \to a\).
\end{solution}

\begin{exercise}
Sinifunktion keskeinen ominaisuus on, että sen arvot kuuluvat välille
\([-1, 1]\) eli \(\sin(x) \in [-1, 1]\) kaikilla \(x \in \RR\). Määritellään
lukujono \((x_n)\) asettamalla, että
\[
x_n = \frac{12\sin(\frac{3\pi n}{8})^4}{n(2 + \sin(\frac{\pi n}{2}))}
\]
kaikilla \(n \in \NNone\). Osoita suoraan määritelmästä, että lukujono \((x_n)\)
suppenee kohti lukua \(0\).
\end{exercise}

\begin{solution}
Olkoon \(\epsilon > 0\). Koska \(\abs{\sin(x)} \leq 1\) kaikilla \(x \in \RR\),
on myös \(\abs{\sin(x)^4} \leq 1\) kaikilla \(x \in \RR\). Täten
\(\abs{12\sin(x)^4} \leq 12\) ja \(\abs{n(2 + \sin(x))} \geq n\).
Valitaan \(N \in \NNone\) siten, että \(N > \frac{12}{\epsilon}\). Tällöin
\[
\abs{x_n - 0} \leq \frac{12}{n} < \frac{12}{N} < \epsilon,
\]
kun \(n > N\). Täten \(\lim_{n \to \infty} x_n = 0\).
\end{solution}
\end{document}
